\documentclass[12pt]{article}
\def\activityid{F06-M-v2}\usepackage[letterpaper,margin=.65in,top=.60in,bottom=.65in]{geometry}
\usepackage[T1]{fontenc}\usepackage[default]{sourcesanspro}
\usepackage{microtype,tikz,array,tabularx,fancyhdr,amsmath}\usepackage[hidelinks]{hyperref}
\usetikzlibrary{arrows.meta,calc}
\setlength{\parindent}{0pt}\setlength{\parskip}{7pt}\setlength{\headheight}{0pt}\setlength{\footskip}{26pt}\setlength{\emergencystretch}{2em}
\pagestyle{fancy}\fancyhf{}\renewcommand{\headrulewidth}{0pt}\renewcommand{\footrulewidth}{.4pt}
\fancyfoot[L]{\scriptsize Bellingham Math Circle / Week 6 / \activityid}\fancyfoot[R]{\scriptsize \thepage}
\definecolor{ink}{gray}{.12}\definecolor{soft}{gray}{.95}\color{ink}
\newcommand{\PageTitle}[2]{{\fontsize{10}{12}\selectfont\bfseries WEEK 6\quad CODE LAB\hfill #1\par}\vspace{3pt}{\fontsize{26}{29}\selectfont\bfseries #2\par}\vspace{2pt}}
\newcommand{\NameLine}{{\small Name: \rule{2.65in}{.4pt}\hfill Date: \rule{1.35in}{.4pt}\par}\vspace{4pt}}
\newcommand{\Task}[2]{\par\vspace{3pt}{\large\bfseries #1}\par #2\par}
\newcommand{\Subhead}[1]{\par\vspace{3pt}{\large\bfseries #1}\par}
\newcommand{\RuleBox}[1]{\par\begingroup\setlength{\fboxsep}{8pt}\colorbox{soft}{\parbox{\dimexpr\linewidth-16pt\relax}{#1}}\endgroup\par}
\newcommand{\WriteBox}[2]{\par\textbf{#1}\par\vspace{2pt}\begin{tikzpicture}\draw[black!40,line width=.5pt] (0,0) rectangle (\dimexpr\linewidth-.5pt\relax,#2);\end{tikzpicture}\par}
\newcommand{\StudentFooter}{\vfill{\small Build first. Explain by pointing, talking, or drawing.}\par}
\newcommand{\Code}[1]{\texttt{#1}}

\hypersetup{pdftitle={Week 6: grades-2-3}}
\begin{document}
\PageTitle{GRADES 2--3}{A method that works for every secret}\NameLine
\RuleBox{Two colors: R and B. Repeats allowed. A test earns one point for each color in the correct position. You hear \textbf{only the total}, not which positions match.}\Task{1. Explore three places.}{Play a round without a test limit. Keep the old tests and scores visible. Then try this investigation.}
\Task{2. Change exactly one place.}{Compare \Code{RRR} with \Code{BRR}. Can the score stay the same? What does an increase or a decrease tell you?}
\begin{center}\renewcommand{\arraystretch}{2}\begin{tabular}{c|p{2.2in}|p{.7in}|p{1.5in}}Test&My code&Score&Still possible\\\hline 1&&&\\2&&&\\3&&&\\4&&&\end{tabular}\end{center}
\Task{3. Three tests, then name the secret.}{Try \Code{RRR}, \Code{BRR}, and \Code{RBR}. How can these reveal all three places? Check your method on a partner's secret.}
\WriteBox{Why my method always determines the third place too:}{1.05in}
\textbf{Counting rule:} Naming a known secret costs no test here. In the original game, submitting it as a winning guess may cost one more turn.\StudentFooter
\newpage
\PageTitle{GRADES 2--3}{Can anyone always do it in two?}\NameLine
\RuleBox{Two colors: R and B. Repeats allowed. A test earns one point for each color in the correct position. You hear \textbf{only the total}, not which positions match.}\Task{4. The difficult first answer.}{Suppose \Code{RRR} scores 2. Exactly these secrets remain:}
\begin{center}\fontsize{22}{26}\selectfont\Code{BRR}\qquad\Code{RBR}\qquad\Code{RRB}\end{center}
\Task{5. Can a last test separate all three?}{Try rows of the table and circle repeated scores. What is the only feature that changes among the three secrets? Explain why a test can tell only whether the secret's B is in one of its B positions. \textbf{Stop calculating rows once you can explain why at most two scores occur.}}\begin{center}\renewcommand{\arraystretch}{1.4}\begin{tabular}{c|c|c|c}Test&BRR&RBR&RRB\\\hline RRR&&&\\RRB&&&\\RBR&&&\\RBB&&&\\BRR&&&\\BRB&&&\\BBR&&&\\BBB&&&\end{tabular}\end{center}
\Task{6. Does a different first test help?}{Try \Code{BRB}. In \emph{both} a secret and every test, swap R/B at places 1 and 3. The first test becomes \Code{RRR}. Check that each matching place still matches. Do this at every B position of any first test. Why are all scores unchanged? Discuss with an adult.}
\WriteBox{Why no first test can guarantee success in only two tests:}{.65in}
\textbf{What we proved:} Task 5 rules out two tests starting with RRR. With task 6, it covers every first test; together with page 1, three is optimal.\StudentFooter
\end{document}
